\chapter{Comparación con el test beam}
El observable experimental se define como la pendiente de $\langle t_L-t_R\rangle$ frente a la posición. El reporte de velocidad efectiva tabula la convención de dos extremos y las convenciones one-end, sin escoger entre ellas.

La reconstrucción CFD de EXEC\_46 usa tres formas de pulso evaluadas sobre la suma temporal de hits: dos modelos con $(\tau_r,\tau_f)=(2,3)\,\mathrm{ns}$ y un contraste rápido $(0.5,2.0)\,\mathrm{ns}$. Ninguno parametriza la cadena MUSIC/SAMPIC del test beam. Por tanto, la comparación con el valor de referencia de Betancourt se etiqueta como indicativa.

Las tablas de `veff\_rank/` contienen pendientes, chi2/ndf, errores, fracciones CFD y theta efectivo. La columna `path\_over\_daxial\_exact` usa, evento a evento, $|x_{SiPM}-x_{creation}|$ con $x_{SiPM}=\pm700\,\mathrm{mm}$. No se generan figuras nuevas para esta sección.

\begin{figure}[ht]
\centering
\includegraphics[width=0.82\textwidth]{v2_g_d.pdf}
\caption{Transporte $g(d)$ de la campaña v2 con óptica dispersiva.}
\end{figure}
